% ============================================================================
% DRAFT of the RQ3 subsection for the Results section — to be rewritten by the authors in their own words.
% Every number below comes from rq3/out/tables/*.csv (saved fits, Level 4). "TBD" = not computed yet;
% all numbers come from the reduced pilot design (20k training applicants, five retrainings) and must stay labelled preliminary.
% \cite keys are placeholders: replace with the keys of your .bib (A2 refs [13] Lundberg & Lee 2017; [21] Lundberg et al. 2020).
% ============================================================================

\subsection{RQ3: Consistency and Stability of SHAP Explanations}
\label{subsec:rq3}

RQ3 asks how consistent SHAP explanations of credit decisions are between the simple and the complex models, and how stable they remain after the models are retrained. H3 states that the complex models give less stable explanations than logistic regression and agree with it only partly on the main reasons for a decision. Following the Methods, two measures are used: the Spearman correlation between the global feature rankings of two explanations (features ranked by mean $|\mathrm{SHAP}|$ over the explained applicants) and the share of an applicant's five largest reasons that are present in both explanations. Explanations are called stable when the mean Spearman correlation is at least 0.8; the poster added the applicant-level rule that at least four of the five top reasons must stay the same, which is reported alongside. Logistic regression is the reference for every comparison. All results of this subsection are preliminary: they are computed on the pilot models of Section~\ref{subsec:rq1} (20,000 training applicants, five search configurations) and on the retraining design described below.


\begin{table}[t]
\caption{Set-up of the explanation experiment (RQ3).}
\label{tab:rq3_setup}
\centering
\small
\begin{tabular}{p{2.6cm}p{5.4cm}}
\hline
Item & Value \\
\hline
Explained applicants & 1,000 from the locked test set (500 with, 500 without bureau record), seed 42; 300 for the ensembles; 500 (300 for the ensembles) in the retraining runs \\
Default rate among them & 8.6\% in both groups \\
Background sample & 100 applicants of the 20k pilot training set, seed 42 \\
Feature level & Level 4, 322 raw features (755 columns after preprocessing) \\
Explained output & probability of default $P(\mathrm{TARGET}=1)$ \\
Exact explainers & interventional Tree SHAP: random forest, XGBoost, LightGBM \\
Approximate explainers & permutation Shapley: LR, MLP (2 permutations on the saved fits, 1 in the retraining runs), soft voting, stacking (1 permutation) \\
Retrainings & 5 per model, seeds 1--5, 90\% stratified subsample of the 20k training set \\
\hline
\end{tabular}
\end{table}

\paragraph{Set-up and deviations from the plan.}
Table~\ref{tab:rq3_setup} summarises the experiment. One list of 1,000 test applicants was fixed once and reused for every model and run: 500 applicants with and 500 without a credit-bureau record were drawn at random (seed 42) from the locked test set, which keeps each group's natural default rate (8.6\% in both groups). The background sample used by every explainer consists of 100 applicants drawn at random (seed 42) from the 20,000 pilot training applicants; no test applicant is used as background. All explanations are computed at history Level~4 (322 raw features; 755 columns after preprocessing) on one common scale, the predicted probability of default $P(\mathrm{TARGET}=1)$. SHAP values of the one-hot and missing-indicator columns are summed back to their raw feature, so all models are compared on the same 322 features. Random forest, XGBoost and LightGBM are explained with the exact interventional Tree SHAP algorithm \cite{lundberg2020trees}; logistic regression, the MLP, soft voting and stacking are explained with a permutation Shapley estimator (forward and reverse pass over a random feature order, averaged over the background sample; two permutations for logistic regression and the MLP, one for the ensembles). For every model on the saved fits, the sum of the SHAP values plus the base value reproduced the model output to within $8\times10^{-5}$. In the seed-5 retraining of LightGBM, the maximum absolute additivity residual was 0.001374, with one of 500 applicants exceeding $10^{-3}$; this is an additivity diagnostic rather than a ranking-stability measure, and the original explanations were retained. Other retrainings also exceeded the saved-fit additivity bound; their raw diagnostics were retained. Three deviations from Assignment~2 apply. First, logistic regression is explained with the permutation estimator on the probability scale instead of the LinearExplainer, because the latter explains log-odds and one common scale was required for all models; on a pilot sample of 10 applicants the exact log-odds ranking of logistic regression and the permutation estimate agreed with a global Spearman correlation of 0.99 and a top-5 overlap of 0.96. Second, to stay within the computation budget, the ensembles were explained on the first 150 applicants of each group (300 in total) instead of 1,000, and the retraining runs explained the first 250 applicants of each group (500; 300 for the ensembles) with one permutation; every table states the number of applicants used. Third, the five retrainings use the saved best hyper-parameters and a stratified 90\% subsample of the 20,000 pilot training applicants (seeds 1--5), not of the full development set; seeds 1--4 were run in the cloud container and seed 5 on a separate laptop described in the reproducibility paragraph, with the same package versions, applicant list and background sample.

\input{rq3_consistency_with_LR}

\begin{figure}[t]
\centering
\includegraphics[width=0.95\columnwidth]{fig_rq3_agreement_with_LR.pdf}
\caption{Share of an applicant's five largest SHAP reasons that coincide with those of logistic regression, by availability of a credit-bureau record (saved fits, Level 4; bars: mean, whiskers: 95\% bootstrap CI over applicants).}
\label{fig:rq3_agree}
\end{figure}

\begin{figure}[t]
\centering
\includegraphics[width=0.85\columnwidth]{fig_rq3_spearman_matrix.pdf}
\caption{Spearman correlation between the global feature rankings (mean $|\mathrm{SHAP}|$ over the explained applicants) of every pair of models, saved fits, Level 4.}
\label{fig:rq3_matrix}
\end{figure}

\begin{figure}[t]
\centering
\includegraphics[width=0.95\columnwidth]{fig_rq3_top_features.pdf}
\caption{Rank of the most important raw features (mean $|\mathrm{SHAP}|$, saved fits) for logistic regression, LightGBM and stacking; features in the top eight of at least one of the three models.}
\label{fig:rq3_top}
\end{figure}

\paragraph{Agreement with logistic regression.}
Table~\ref{tab:rq3_agree} reports the agreement of every model with logistic regression on the saved fits. On average, between 1.4 and 1.8 of an applicant's five largest reasons coincide with those of logistic regression: random forest 0.28 [0.27, 0.29], XGBoost 0.37 [0.36, 0.38], LightGBM 0.35 [0.34, 0.36] and the MLP 0.32 [0.31, 0.33] ($n$ = 1,000); the two ensembles, whose base learners include logistic regression itself, share more: soft voting 0.54 [0.52, 0.56] and stacking 0.68 [0.66, 0.70] ($n$ = 300), that is 2.7 and 3.4 of five reasons. The Spearman correlation between the global rankings of a complex model and of logistic regression lies between 0.54 and 0.58 for the tree models, reaches 0.66 for the MLP, 0.81 [0.81, 0.82] for soft voting and 0.87 [0.87, 0.88] for stacking. The tree models agree much more with one another than with logistic regression: the global rankings of XGBoost and LightGBM correlate at 0.93, those of random forest with either boosting model at 0.84--0.86 (Fig.~\ref{fig:rq3_matrix}), and XGBoost and LightGBM share 0.73 of their top-5 reasons per applicant. Agreement with logistic regression is lower for applicants without a bureau record than for applicants with one for random forest (0.24 versus 0.32) and slightly lower for XGBoost (0.36 versus 0.37) and LightGBM (0.33 versus 0.36), whereas the MLP shows no difference (Fig.~\ref{fig:rq3_agree}). The two external scores EXT\_SOURCE\_2 and EXT\_SOURCE\_3 are the two most important features for logistic regression, random forest, XGBoost and LightGBM alike; from the third rank onward the lists diverge (Fig.~\ref{fig:rq3_top}): logistic regression continues with the average number of future instalments, the annuity and the maximum payment count of previous applications, while the tree models continue with the credit-to-goods ratio, EXT\_SOURCE\_1 and the share of late instalments. The MLP ranks the indicators HAS\_CC and HAS\_BB, which only record whether credit-card and bureau-balance history exists, above every external score.

\IfFileExists{rq3_stability.tex}{\input{rq3_stability}}{}
\IfFileExists{rq3_refit_test_metrics.tex}{\input{rq3_refit_test_metrics}}{}
\IfFileExists{fig_rq3_stability.pdf}{%
\begin{figure*}[t]
\centering
\includegraphics[width=0.9\textwidth]{fig_rq3_stability.pdf}
\caption{Stability of SHAP explanations across retrainings of each model (Level 4): (a) Spearman correlation between the global feature rankings of two retrainings (mean over run pairs, 95\% bootstrap CI over applicants); (b) distribution over applicants of the share of top-5 reasons kept between two retrainings. Dashed line: pre-set criterion 0.8.}
\label{fig:rq3_stability}
\end{figure*}}{}

\paragraph{Stability after retraining.}
Table~\ref{tab:rq3_stability} and Fig.~\ref{fig:rq3_stability} report the stability of each model's explanations across its five retrainings (ten run pairs). The global feature rankings are reproducible for every model except the MLP: the mean Spearman correlation between two retrainings is 0.89 $\pm$ 0.02 for logistic regression, 0.97 $\pm$ 0.00 for random forest, 0.93 $\pm$ 0.01 for XGBoost, 0.89 $\pm$ 0.01 for LightGBM, 0.94 $\pm$ 0.01 for soft voting and 0.89 $\pm$ 0.02 for stacking, all above the pre-set threshold of 0.8, whereas the MLP reaches 0.79 $\pm$ 0.01 [0.76, 0.80]. At the applicant level only random forest satisfies the four-of-five rule: it keeps 0.81 [0.80, 0.82] of an applicant's five largest reasons between two retrainings, logistic regression keeps 0.78 [0.77, 0.79], soft voting 0.77 [0.76, 0.78], LightGBM 0.74 [0.74, 0.75], stacking 0.74 [0.73, 0.75], XGBoost 0.69 [0.68, 0.70] and the MLP 0.30 [0.29, 0.31]. The paired differences in applicant-level stability relative to logistic regression, with a Bonferroni-adjusted interval for the six comparisons, are +0.03 [+0.01, +0.05] for random forest, $-$0.09 [$-$0.11, $-$0.08] for XGBoost, $-$0.04 [$-$0.06, $-$0.02] for LightGBM, $-$0.48 [$-$0.50, $-$0.46] for the MLP and $-$0.04 [$-$0.06, $-$0.02] for stacking, all excluding zero, while soft voting ($-$0.01 [$-$0.03, +0.01]) does not differ from logistic regression. Under the pre-set criteria, therefore, H3 is supported for XGBoost, LightGBM, stacking and the MLP, not supported for soft voting, and reversed for random forest, whose applicant-level explanations are more stable than those of logistic regression. The predicted probabilities themselves move little between retrainings: the mean standard deviation of an applicant's predicted default probability over the five runs is 0.007--0.015 for all models except the MLP (0.080). Table~\ref{tab:rq3_refit_metrics} gives the test-set metrics of the retrained models on the full locked test set as mean $\pm$ s.d.\ over the five retrainings; the standard deviations of PR-AUC lie between 0.0009 (LightGBM) and 0.0020 (stacking) and those of ROC-AUC at most 0.0037 (MLP), which supplies the retraining variability that the single fits of Section~\ref{subsec:rq1} could not report.

\paragraph{Explainer noise and sub-groups.}
For the permutation-based explanations, the saved fits were explained twice with independent random feature orders at the settings of the retraining runs (one permutation). Two such explanations of the same logistic-regression fit share 0.90 [0.89, 0.91] of an applicant's top-5 reasons ($n$ = 500), of the same soft-voting fit about 0.88 [0.86, 0.89] and of the same stacking fit 0.77 [0.75, 0.80] ($n$ = 150), while the global rankings agree at a Spearman correlation of 1.00 for all three; the exact Tree SHAP explanations are deterministic. The applicant-level stability of stacking (0.74) is therefore close to the noise floor of its explainer, and that of soft voting (0.77) and logistic regression (0.78) lies about 0.11 below their floors, so same-fit noise provides a conditional reference and cannot isolate the causal contribution of refitting; the tree models have no such floor. For the MLP the floor itself is 0.39 [0.37, 0.41], confirming that its applicant-level explanations are not reproducible even without retraining. Stability is similar with and without a credit-bureau record: the share of top-5 reasons kept is 0.80 versus 0.76 for logistic regression, 0.80 versus 0.81 for random forest, 0.68 versus 0.69 for XGBoost, 0.74 versus 0.74 for LightGBM, 0.78 versus 0.76 for soft voting and 0.76 versus 0.72 for stacking; the global Spearman correlation of the MLP is 0.80 with and 0.74 without a bureau record, the only case in which one group falls below the threshold.

\paragraph{Error analysis and negative results.}
The MLP gives unfavourable explanation diagnostics. On the saved Level-4 fit, 38.9\% of the explained applicants receive a predicted default probability below $10^{-6}$ and 0.5\% a probability above $1-10^{-6}$, so many applicants receive extreme probabilities; the median sum of absolute SHAP values per applicant is 1.05 for the MLP against 0.59 for logistic regression and 0.14--0.20 for the tree models, and the pilot noise check gave a top-5 overlap of only 0.32 between two explanations of the same MLP fit. The MLP is kept in all tables, as planned, but its explanation measures must be read together with this saturation. Across the five retrainings EXT\_SOURCE\_2 stays the first-ranked feature of logistic regression, LightGBM and stacking in every run and EXT\_SOURCE\_3 stays second for LightGBM and stacking, whereas the ranks below the top two move: for logistic regression PREVIOUS\_AMT\_CREDIT\_MAX ranges from rank 4 to 21, FLAG\_DOCUMENT\_3 from 6 to 16 and INSTALLMENTS\_RECORD\_COUNT from 4 to 10, for LightGBM EXT\_SOURCE\_1, INSTALLMENTS\_LATE\_SHARE and APP\_CREDIT\_TO\_GOODS exchange ranks 3 to 5, and for stacking INSTALLMENTS\_DUE\_DAY\_MAX ranges from 7 to 19 and POS\_CNT\_INSTALMENT\_FUTURE\_AVG from 9 to 17. The share of applicants who keep fewer than three of their five reasons between two retrainings is 8.2\% for logistic regression, 2.2\% for random forest, 19.6\% for XGBoost, 5.6\% for LightGBM, 5.3\% for soft voting, 9.0\% for stacking and 98.8\% for the MLP. For random forest and soft voting these applicants have a higher mean predicted default probability than the stable ones (0.10 and 0.14 versus 0.08). Predicted-risk differences are descriptive and carry no causal interpretation.

% ---- Reproducibility paragraph for the whole Results section (RQ3 part; merge with Askat's) -------------------
\paragraph{Reproducibility.}
The code of all three questions is available at \url{https://github.com/AskatQwerty/RMTpreresults} (commit \texttt{4df488e}, branch \texttt{pilot-handoff}; the RQ3 scripts are in \texttt{rq3/}). Models were trained on a laptop with an AMD Ryzen~7 7435HS (8 cores, 3.1~GHz) and 16~GB RAM under Windows~11 with Python~3.13.9, scikit-learn~1.7.2, XGBoost~3.2.0 and LightGBM~4.7.0; the whole pilot search and refit took 10,640~s. The saved-fit explanations and retraining seeds 1--4 ran in a Linux container with an Intel Xeon processor (2 cores, 2.8~GHz) and 7~GB RAM under Python~3.13.17 with the same library versions and shap~0.54.1. Retraining seed 5 ran on a separate laptop with an Intel Core i5-10300H processor (4 cores, 8 threads), 8~GB RAM and Windows under Python~3.13.1, using the same pinned library versions, explained applicants and background sample. Seeds: 42 for the train/test split, the search and all samples; 1--5 for the retrainings; 0, 11 and 12 for the permutation explainer. Explaining the saved fits took about 35~min in total (exact Tree SHAP: under 3~min for the three tree models; permutation estimator: 0.2--1.4~s per applicant depending on the model); one retraining of all seven models took about 25~min in the container, of which the random forest fit alone took 7.5~min and the other six fits between 4~s and 80~s.
